\section{Numerical Calculus}
\label{sec:numcalc}

Every operation in the previous sections was \emph{exact}: \B{Integrate} returns an
antiderivative, \B{Solve} returns the roots as algebraic numbers, \B{Sum} returns a closed
form. That is the ideal, and when it is attainable it is unbeatable. But most of the
integrals, equations, sums, and differential equations that arise in practice have no closed
form at all---not because Mathilda is not clever enough, but because, as theorems of Liouville,
Abel, and others tell us, none exists. For those problems the honest answer is a
\emph{number}, computed to whatever precision you ask for, and computing it well is a
subject in its own right.

This section is about Mathilda's \emph{numerical} calculus: the family of functions whose
names begin with \texttt{N}. We start with the question of when to reach for them at all, and
with the one contract they all share---how they decide they are ``done.'' Then we take them in
turn: numerical integration (\B{NIntegrate}), summation and products (\B{NSum}, \B{NProduct}),
differentiation (\B{ND}), series and limits (\B{NSeries}, \B{NLimit}), residues (\B{NResidue}),
differential equations (\B{NDSolve}), the roots of polynomials and the solutions of general
equations (\B{NRoots}, \B{NSolve}, \B{FindRoot}), and finally optimization, both local
(\B{FindMinimum}, \B{FindMaximum}) and global (\B{NMinimize}, \B{NMaximize}). For each we show
not just the syntax but the algorithm underneath---because choosing the right method for the
shape of your problem is most of the craft of numerical work.

\subsection{Why and when to go numeric}
\label{ssec:whynumeric}

The clearest way to see the need is to watch a symbolic function stop short and a numerical one
carry on. The Gaussian \(\int_0^1 e^{-x^2}\,dx\) has no elementary antiderivative; \B{Integrate}
does the best exact thing it can, expressing the answer through the error function, while
\B{NIntegrate} simply returns the number:

\begin{codepairs}
\pair{numerical-calculus/motivation/1}{No elementary antiderivative exists, so the exact answer
is written through a special function, \(\tfrac12\sqrt{\pi}\,\operatorname{erf}(1)\).}
\pair{numerical-calculus/motivation/2}{\B{NIntegrate} gives the value you usually wanted in the
first place.}
\pair{numerical-calculus/motivation/3}{The sine integral is worse: it has no antiderivative
elementary in \(x\) at all, so \B{Integrate} names it \(\operatorname{Si}(1)\).}
\pair{numerical-calculus/motivation/4}{Again the number is one step away.}
\end{codepairs}

Roots tell the same story. Abel proved that the general quintic has no solution in radicals, so
\B{Solve} can only name the five roots of \(x^5-x-1\) as abstract \texttt{Root} objects---correct,
but not a number you can plot---whereas \B{NRoots} and \B{FindRoot} deliver digits:

\begin{codepairs}
\pair{numerical-calculus/motivation/5}{There is no radical formula, so each root is returned as
an indexed \texttt{Root} object standing for ``the \(k\)th root of this polynomial.''}
\pair{numerical-calculus/motivation/6}{\B{NRoots} approximates all five at once: one real, two
conjugate pairs.}
\pair{numerical-calculus/motivation/7}{\B{FindRoot} chases a single root of a transcendental
equation---\(\cos x = x\), whose solution (the ``Dottie number'') is provably not expressible in
closed form.}
\end{codepairs}

The rule of thumb: reach for a numerical function when the exact machinery returns something
unevaluated, something written through a special function or \texttt{Root} object you would only
numericalize anyway, or something correct but too large to be useful---and whenever the input is
itself inexact, since there is then no exact answer to find.

\subsection{Precision, and the three goals}
\label{ssec:numgoals}

Every numerical function in this section shares one control panel, so it is worth learning once.
Three options govern how hard the function works and how it judges success:
\texttt{WorkingPrecision}, \B{AccuracyGoal}, and \B{PrecisionGoal}\index{numerical calculus!accuracy
and precision goals}.

\texttt{WorkingPrecision} sets how many digits the internal arithmetic carries. Its default is
\texttt{MachinePrecision}---the roughly \(15.95\) decimal digits of an IEEE double
(\S\ref{ssec:machine-precision})---and a larger value transparently switches the computation onto
Mathilda's MPFR path, so you can ask for forty digits or five hundred:

\begin{codepairs}
\pair{numerical-calculus/goals/1}{At machine precision the answer is a plain double.}
\pair{numerical-calculus/goals/2}{Raising \texttt{WorkingPrecision} to \(40\) recomputes the same
integral in extended precision.}
\pair{numerical-calculus/goals/3}{The same lever works for a sum: forty digits of \(\pi^2/6\).}
\end{codepairs}

The other two options set the \emph{target}, not the arithmetic. \B{AccuracyGoal}~\(a\) asks for a
result whose \emph{absolute} error is below \(10^{-a}\); \B{PrecisionGoal}~\(p\) asks for a
\emph{relative} error below that fraction of the answer's size. A function stops when it has met the
combined tolerance
\[
  \text{error} \;\le\; 10^{-a} \;+\; |x|\,10^{-p},
\]
where \(x\) is the current estimate.\calloutnote{Under the hood}{One small module,
\texttt{src/nc\_accuracy.c}, implements this contract for every routine here, so the semantics never
drift between them. \texttt{Automatic} resolves to two digits below \texttt{WorkingPrecision}; the symbol
\texttt{MachinePrecision} tracks the machine-double budget; and \texttt{Infinity} switches a
criterion off, so \texttt{PrecisionGoal -> Infinity} judges by absolute error alone. The near-universal
defaults are \texttt{AccuracyGoal -> MachinePrecision} and \texttt{PrecisionGoal -> Automatic}.}
The absolute term keeps the relative term from demanding the impossible near a zero answer, and the
relative term keeps the absolute term honest for very large or very small ones.

One behavior is worth internalizing now, because you will meet it throughout the section:
\textbf{these functions never fail with an error.} If a routine cannot reach its goal within its
resource budget, it refines as far as it can, keeps the best estimate it found, prints a
\texttt{Head::accgl} message to the terminal, and returns that estimate
anyway.\calloutnote{Pitfall}{The \texttt{::accgl} message is advisory, not fatal, and it is printed
to standard error---so it appears at your prompt but not inside a captured result. When a numerical
answer looks a digit or two short, it is often because the goal was quietly not met; raise
\texttt{WorkingPrecision}, loosen the goal, or (as we will see) pick a method better suited to the
problem.} We rely on this deliberately below, where forcing the wrong summation method leaves a
visibly wrong last digit.

\subsection{Numerical integration}
\label{ssec:nintegrate}

\B{NIntegrate}\index{NIntegrate@\mcode{NIntegrate}!methods} is the workhorse, and its default
\texttt{Automatic} method is really a small polyalgorithm that inspects the integrand and the
interval and dispatches to a specialist. On a smooth, finite integrand it uses globally-adaptive
Gauss--Kronrod quadrature; on a singular or infinite one it switches to a double-exponential rule;
on an oscillatory one it uses a Levin collocation scheme; and in high dimensions it falls back to
Monte-Carlo sampling. The everyday cases just work:

\begin{codepairs}
\pair{numerical-calculus/nintegrate-basic/1}{An infinite interval: \(\int_0^\infty e^{-x^2}\,dx =
\tfrac12\sqrt{\pi}\).}
\pair{numerical-calculus/nintegrate-basic/2}{An oscillatory tail: the Dirichlet integral
\(\int_0^\infty \tfrac{\sin x}{x}\,dx = \tfrac{\pi}{2}\).}
\pair{numerical-calculus/nintegrate-basic/3}{An integrable endpoint singularity at \(x=0\); the
adaptive rule clusters its samples there.}
\pair{numerical-calculus/nintegrate-basic/4}{Two dimensions: adaptive cubature over the plane gives
\(\int_{\mathbb{R}^2} e^{-x^2-y^2} = \pi\).}
\end{codepairs}

\usagebox{NIntegrate}

When you know the shape of your integrand, naming the \texttt{Method} is both faster and more
accurate. The most useful members of the family are \texttt{"GlobalAdaptive"} (the Gauss--Kronrod
default)\index{Gauss--Kronrod quadrature}, \texttt{"DoubleExponential"} (the tanh-sinh / sinh-sinh /
exp-sinh rules, superb for endpoint singularities and infinite ranges)\index{double-exponential
quadrature}, \texttt{"LevinRule"} (for products with an oscillatory kernel like \(\cos g\), \(\sin
g\), or \(e^{ig}\))\index{Levin collocation}, \texttt{"MonteCarlo"} and its quasi-random variant (for
many dimensions and for region indicators), and \texttt{"PrincipalValue"} (for a Cauchy principal
value around a pole named in \texttt{Exclusions}):

\begin{codepairs}
\pair{numerical-calculus/nintegrate-methods/1}{The double-exponential rule reproduces the Gaussian on
the half-line to machine precision.}
\pair{numerical-calculus/nintegrate-methods/2}{A rapidly oscillating kernel: Levin collocation returns
\(\int_0^\infty \cos(100x)\,e^{-x}\,dx = 1/10001\) without resolving every wiggle.}
\pair{numerical-calculus/nintegrate-methods/3}{A Cauchy principal value straddling the pole at the
origin: the divergent halves cancel to \(\ln 2\).}
\pair{numerical-calculus/nintegrate-methods/4}{Monte-Carlo sampling of a two-dimensional Gaussian; the
seed is fixed, so the estimate is reproducible.}
\end{codepairs}

\calloutnote{Theory}{Gauss--Kronrod quadrature\index{Gauss--Kronrod quadrature} is the engine behind
the default. A Gauss rule of \(n\) points integrates polynomials up to degree \(2n-1\) exactly; the
Kronrod extension adds \(n+1\) points so that the \(15\)-point rule reuses every node of the embedded
\(7\)-point rule. The difference between the two estimates is a cheap, reliable error bar, and the
adaptive driver bisects whichever subinterval carries the largest one until the total error meets the
goal---the classic QUADPACK strategy~\cite{quadpack}.}

\calloutnote{Theory}{The double-exponential (tanh-sinh) rule~\cite{takahasimori} is a change of variable
\(x = \tanh\!\big(\tfrac{\pi}{2}\sinh t\big)\) that makes the integrand decay doubly exponentially at
the ends. A plain trapezoidal rule in \(t\) then converges with astonishing speed, and---crucially---it
places no sample exactly at the endpoints, so an integrable singularity there costs nothing. This is why
it is Mathilda's method of choice for singular and infinite integrals and for high precision.}

Two honesty notes. Several rule names the reference documentation lists---\texttt{"ClenshawCurtisRule"},
\texttt{"LobattoKronrodRule"}, and a few others---are recognized but not yet implemented, and asking for
one yields an \texttt{NIntegrate::method} message rather than a silent substitution. And
\texttt{"AdaptiveMonteCarlo"} currently runs the same sampler as \texttt{"MonteCarlo"} without the
adaptive stratification its name promises. The book prefers to say so plainly.

\subsection{Numerical summation and products}
\label{ssec:nsum}

An infinite sum is an integral's discrete cousin, and \B{NSum}\index{NSum@\mcode{NSum}!methods} carries
the same idea: add the first several terms explicitly, then \emph{extrapolate} the tail rather than
plod through millions of terms. Its \texttt{Automatic} method reads the series and chooses---the
Euler--Maclaurin formula for a smooth monotone tail\index{Euler--Maclaurin summation}, the
Cohen--Villegas--Zagier algorithm for an alternating one\index{alternating series!acceleration}, and
Wynn's epsilon algorithm otherwise\index{Wynn epsilon algorithm}:

\begin{codepairs}
\pair{numerical-calculus/nsum/1}{The Basel sum \(\sum 1/n^2 = \pi^2/6\); Euler--Maclaurin nails it from
a handful of terms.}
\pair{numerical-calculus/nsum/2}{An alternating series for \(\pi/4\); the \texttt{"AlternatingSigns"}
method converges geometrically where naive addition crawls.}
\pair{numerical-calculus/nsum/3}{A slowly converging \(p\)-series at forty digits---compare
\(\zeta(11/10)\). Adding terms directly, you would need more than \(10^{400}\) of them.}
\end{codepairs}

\usagebox{NSum}

Extrapolation is powerful but not free, and the wrong accelerator can hurt. Forcing Wynn's epsilon
onto the Basel sum---whose tail is monotone, not oscillatory---leaves the last digits wrong:

\begin{codepairs}
\pair{numerical-calculus/nsum/4}{Compare with pair~1 above: \(1.62533\) instead of \(1.64493\). Wynn's
epsilon is the wrong tool for a monotone tail, and (because it cannot reach the goal) an
\texttt{NSum::accgl} message accompanies this at the terminal.}
\end{codepairs}

\calloutnote{Theory}{Wynn's epsilon algorithm\index{Wynn epsilon algorithm} is a recursive
implementation of the Shanks transformation, which models a sequence's tail as a sum of decaying
geometric terms and cancels them. It is spectacular on alternating and geometric-like sequences and
weak on the slow algebraic decay of a monotone \(p\)-series---exactly the failure the pair above shows.
The Euler--Maclaurin formula, by contrast, replaces the monotone tail with an integral plus a
correction series in the Bernoulli numbers, which is why \texttt{Automatic} prefers it there.}

\B{NProduct}\index{NProduct@\mcode{NProduct}} handles infinite products by the obvious device---it sums
the logarithms and exponentiates, so \(\prod f = \exp\!\sum \log f\)---and therefore inherits \B{NSum}'s
methods and convergence test wholesale:

\begin{codepairs}
\pair{numerical-calculus/nproduct/1}{\(\prod_{n\ge 2}\!\big(1-\tfrac1{n^2}\big) = \tfrac12\), a
telescoping classic.}
\pair{numerical-calculus/nproduct/2}{\(\prod_{n\ge 2}\tfrac{n^2}{n^2-1} = 2\), its reciprocal.}
\end{codepairs}

\usagebox{NProduct}

\subsection{Numerical differentiation}
\label{ssec:nd}

Differentiation is exact and easy when you have a formula, but a numerical derivative is what you need
when the function is only known through a black box, or is not analytic, or when a symbolic derivative
would be enormous. \B{ND}\index{ND@\mcode{ND}!methods} estimates the derivative of \texttt{expr} at a
point. Its default \texttt{EulerSum} method samples finite differences on a shrinking step and
Richardson-extrapolates them to the limit\index{Richardson extrapolation}:

\begin{codepairs}
\pair{numerical-calculus/nd/1}{\(\tfrac{d}{dx}e^x\) at \(x=1\) is \(e\).}
\pair{numerical-calculus/nd/2}{A high-order, high-precision derivative: the first derivative of
\(\sin(x^2)\) at \(x=3\) to forty digits.}
\pair{numerical-calculus/nd/3}{The integrand \(\operatorname{Re}\,\cos(iy) = \cosh y\) is not
analytic in the complex sense, so only the finite-difference method applies; its derivative at
\(1\) is \(\sinh 1 = 1.1752\).}
\end{codepairs}

\usagebox{ND}

\calloutnote{Under the hood}{\B{ND}'s second method, \texttt{Method -> NIntegrate}, evaluates the
derivative through Cauchy's integral formula, \(f^{(n)}(x_0) = \tfrac{n!}{2\pi i}\oint
\tfrac{f(z)}{(z-x_0)^{n+1}}\,dz\), by calling \B{NResidue} on a small circle. That route needs the
function to be analytic near \(x_0\)---but in return it delivers derivatives of \emph{fractional} and
even complex order, which finite differences cannot. The finite-difference \texttt{EulerSum} method is
the default precisely because it makes no analyticity demand, as pair~3 requires.}

\subsection{Numerical series}
\label{ssec:nseries}

\B{NSeries}\index{NSeries@\mcode{NSeries}} produces a whole Taylor or Laurent expansion numerically:
it samples the function on a circle in the complex plane and recovers the coefficients by a
discrete Fourier transform---Cauchy's integral formula, one contour, all coefficients at once. It is
the natural extension of the contour derivative behind \B{ND}, from a single coefficient to the whole
expansion. Its single algorithm has no \texttt{Method} to choose; you control only the sampling
\texttt{Radius}. Because the DFT carries a little round-off into the coefficients it does not use, it is
idiomatic to \B{Chop} the result:

\begin{codepairs}
\pair{numerical-calculus/nseries/1}{The Taylor series of \(e^x\) about \(0\); the coefficients
\(1,1,\tfrac12,\tfrac16,\tfrac1{24}\) are read straight off.}
\pair{numerical-calculus/nseries/2}{A genuine Laurent series: \(\cot x = \tfrac1x - \tfrac{x}{3} -
\tfrac{x^3}{45} - \cdots\), pole and all.}
\end{codepairs}

\usagebox{NSeries}

\subsection{Numerical limits}
\label{ssec:nlimit}

\B{NLimit}\index{NLimit@\mcode{NLimit}!methods} evaluates a limit by sampling the function along a
sequence of points approaching the target and accelerating that sequence to its limit---the same
extrapolation family as \B{NSum}, since a limit is what a sequence of partial results is converging to.
It is the numerical companion to the symbolic \B{Limit} of \S\ref{sec:calculus}, and it shines when a
limit is delicate or defined only through a hard special function:

\begin{codepairs}
\pair{numerical-calculus/nlimit/1}{A \(0/0\) form: \(\lim_{x\to 0}(1-\cos x)/x^2 = \tfrac12\). Here
\texttt{AccuracyGoal -> 10} simply asks for ten clean digits.}
\pair{numerical-calculus/nlimit/2}{Two diverging pieces whose difference is finite:
\(\lim_{s\to1}\big(\zeta(s)-\tfrac1{s-1}\big) = \gamma\), the Euler--Mascheroni constant.}
\pair{numerical-calculus/nlimit/3}{An \(\infty\)-point limit computed with Wynn's epsilon
(\texttt{"SequenceLimit"}): \(\lim_{n\to\infty} n\,(2^{1/n}-1) = \ln 2\).}
\end{codepairs}

\usagebox{NLimit}

\subsection{Numerical residues}
\label{ssec:nresidue}

The contour trick that powers \B{NSeries} is a tool in its own right. \B{NResidue}\index{NResidue@%
\mcode{NResidue}} computes the residue of a function at a point---the coefficient of \((z-z_0)^{-1}\) in
its Laurent expansion---by integrating around a small circle with the periodic trapezoidal rule, which
is exponentially accurate on a smooth closed contour. Its reach is exactly the symbolic \B{Residue}'s
blind spot: essential singularities, where no Laurent series terminates.

\begin{codepairs}
\pair{numerical-calculus/nresidue/1}{The simplest pole: \(\operatorname{Res}_{z=0}\tfrac1z = 1\).}
\pair{numerical-calculus/nresidue/2}{A pole of \(\tan z\) at \(\pi/2\) has residue \(-1\).}
\pair{numerical-calculus/nresidue/3}{An \emph{essential} singularity: \(\operatorname{Res}_{z=0}
e^{1/z} = 1\). Here the contour \texttt{Radius} is widened to \(1\)---see the pitfall below.}
\pair{numerical-calculus/nresidue/4}{Cauchy coefficient extraction: the residue of \(e^x/x^4\) at
\(0\) is the \(x^3\) Taylor coefficient of \(e^x\), namely \(1/3! = 0.166667\).}
\end{codepairs}

\usagebox{NResidue}

\calloutnote{Pitfall}{The default contour radius is \(1/100\), which is ideal for an ordinary pole but
dangerous for an essential singularity: \(e^{1/z}\) reaches \(e^{100}\) on a circle of radius \(1/100\),
overwhelming the arithmetic. Widening \texttt{Radius} to \(1\), as pair~3 does, keeps the integrand in range.
The lesson generalizes---when a residue near an essential singularity comes back as noise, enlarge the
contour.}

\subsection{Differential equations}
\label{ssec:ndsolve}

Most differential equations have no closed-form solution, which makes \B{NDSolve}\index{NDSolve@%
\mcode{NDSolve}!methods} one of the most important functions in the system. It integrates an initial-value
problem---a single ODE, a coupled system, a higher-order equation reduced automatically to first order,
or a partial differential equation over a rectangle---and returns the solution not as a formula but as an
\B{InterpolatingFunction}\index{InterpolatingFunction@\mcode{InterpolatingFunction}}: a callable object
you can evaluate, differentiate, or plot anywhere in the solved interval.

\begin{codepairs}
\pair{numerical-calculus/ndsolve-numeric/1}{A \emph{stiff} problem, solved with the implicit
\texttt{"BackwardEuler"} method; the answer is an \texttt{InterpolatingFunction} over \([0,3]\).}
\pair{numerical-calculus/ndsolve-numeric/2}{The harmonic oscillator \(y''+y=0\), \(y(0)=1\),
\(y'(0)=0\); assign the solution to reuse it.}
\pair{numerical-calculus/ndsolve-numeric/3}{Evaluating the solution at \(x=2\)\dots}
\pair{numerical-calculus/ndsolve-numeric/4}{\dots reproduces \(\cos 2\) to full machine precision.}
\pair{numerical-calculus/ndsolve-numeric/5}{A coupled first-order system returns one
\texttt{InterpolatingFunction} per dependent variable.}
\end{codepairs}

\usagebox{NDSolve}

The \texttt{Method} option chooses the integrator. The default, \texttt{"DOPRI5"}, is the adaptive
Dormand--Prince \(5(4)\) Runge--Kutta pair\index{Runge--Kutta method!Dormand--Prince}; for a
\emph{stiff} problem---one where the explicit step size is throttled not by accuracy but by stability---you
want an implicit method such as \texttt{"BackwardEuler"} or the variable-order \texttt{"BDF"}. A nonlinear
oscillator makes the interpolating solution vivid. The van der Pol equation \(y'' - \mu(1-y^2)y' + y = 0\)
relaxes onto a \emph{limit cycle}, a closed orbit no linear equation can produce; plotting \(y'\) against
\(y\) draws it directly:

\plotcell{numerical-calculus/vanderpol}

The trajectory starts near \((1,0)\), spirals outward, and locks onto the sharp-cornered limit cycle
characteristic of a relaxation oscillator.

\calloutnote{Theory}{An explicit Runge--Kutta method advances the solution by sampling the slope at
several trial points inside the step and combining them; the Dormand--Prince pair computes two estimates
of different orders from the \emph{same} seven slope evaluations, and their difference drives the adaptive
step size---large where the solution is placid, small where it turns sharply. A \emph{stiff} equation
defeats this: stability, not accuracy, forces a tiny step, and an \emph{implicit} method---one that solves
an algebraic equation for the next value---removes that ceiling. This explicit/implicit split is the
central fact of practical ODE solving~\cite{hairer}.}

For a partial differential equation, \B{NDSolve} uses the \emph{method of lines}: it discretizes the
spatial derivatives on a grid, turning the PDE into a large system of ODEs in time, and integrates that.
The heat equation \(u_t = u_{xx}\) started from a half-sine profile shows the fundamental mode decaying
exponentially:

\plotcell{numerical-calculus/heat-profile}

\calloutnote{Under the hood}{The spatial grid, its resolution (\texttt{"MinPoints"}), and the finite-%
difference accuracy order (\texttt{"DifferenceOrder"}) are sub-options of \texttt{Method -> \{"MethodOfLines",
...\}}. Mathilda builds the spatial derivative stencils by Fornberg's algorithm, compiles the resulting
right-hand side for speed, and---because diffusion is stiff---integrates the semi-discrete system with the
BDF method by default. The returned object is then a two-dimensional \texttt{InterpolatingFunction} in
\(t\) and \(x\).}

\subsection{Roots of polynomials}
\label{ssec:nroots}

For the special case of a single polynomial equation, \B{NRoots}\index{NRoots@\mcode{NRoots}!methods}
finds \emph{all} the roots at once and returns them as a disjunction. This is where Abel's theorem bites
hardest and numerics earn their keep:

\begin{codepairs}
\pair{numerical-calculus/nroots/1}{A cubic: one real root and a conjugate pair.}
\pair{numerical-calculus/nroots/2}{The quintic \(x^5-x-1\)---no radical form---resolved to five roots.}
\pair{numerical-calculus/nroots/3}{Complex coefficients are fine: the two square roots of \(3+4i\) are
\(\pm(2+i)\).}
\pair{numerical-calculus/nroots/4}{A triple root is returned with its multiplicity, three copies of
\(1\).}
\pair{numerical-calculus/nroots/5}{Wilkinson's notoriously ill-conditioned polynomial \(\prod_{k=1}^{15}
(x-k)\): Mathilda recovers all fifteen integer roots cleanly.}
\pair{numerical-calculus/nroots/6}{Precision is set by \B{PrecisionGoal}: thirty digits of \(\sqrt2\).}
\end{codepairs}

\usagebox{NRoots}

\calloutnote{Under the hood}{At machine precision \B{NRoots}'s \texttt{Automatic} method forms the
companion matrix and finds its eigenvalues with LAPACK---exactly what \texttt{numpy.roots} does---then
polishes each root with one Newton step, which is what keeps Wilkinson's polynomial (pair~5) honest. At
higher precision it switches to the Aberth--Ehrlich method\index{Aberth--Ehrlich method}, which refines
\emph{all} the roots simultaneously, each Newton correction repelled from the others so they cannot
collide~\cite{aberth}. A note on options: unlike its siblings, \B{NRoots} takes no \texttt{WorkingPrecision}---%
you steer its precision with \B{PrecisionGoal} alone.}

\subsection{Equations and systems}
\label{ssec:nsolve}

\B{NSolve}\index{NSolve@\mcode{NSolve}!methods} is the general numerical equation solver. For a single
polynomial it defers to \B{NRoots}; for a polynomial \emph{system} it runs a Gröbner-basis elimination and
reads the solutions off an eigenvalue problem; and for a transcendental equation it seeds \B{FindRoot} from
a grid of starting points. A domain argument restricts the answers:

\begin{codepairs}
\pair{numerical-calculus/nsolve/1}{Over the default complexes, \(x^4=1\) has four roots.}
\pair{numerical-calculus/nsolve/2}{Restricting to \texttt{Reals} keeps only \(\pm1\).}
\pair{numerical-calculus/nsolve/3}{A genuine system: \(x^2+y^2=1\), \(x^3-y^3=2\) has six complex
solutions, found by the Gröbner/eigenvalue engine.}
\pair{numerical-calculus/nsolve/4}{A transcendental equation, \(e^x-x=7\), solved over the reals by
grid-seeded \B{FindRoot}: two roots.}
\end{codepairs}

\usagebox{NSolve}

\calloutnote{Theory}{For a zero-dimensional polynomial system Mathilda computes a Gröbner basis
(\S\ref{sec:algebra}), which makes the quotient ring \(\mathbb{Q}[x]/I\) a finite-dimensional vector space.
Multiplication by each variable is then a linear map on that space, and the solutions are read from the
eigenvalues and eigenvectors of those multiplication matrices---the Möller--Stetter method~\cite{clo}. It
is a beautiful bridge: the algebra of \S\ref{sec:algebra} reduces the geometry of a solution set to the
linear algebra of \S\ref{sec:linalg}.}

When you want \emph{one} root and can supply a starting point, \B{FindRoot}\index{FindRoot@%
\mcode{FindRoot}!methods} is the direct tool. Its behavior follows the starting specification: a single
start uses Newton's method, two starts use the secant method, and a bracket \(\{x, x_0, x_{\min},
x_{\max}\}\) uses Brent's method. It also solves systems, with a symbolic Jacobian:

\begin{codepairs}
\pair{numerical-calculus/findroot/1}{Newton's method from \(x=1\) on \(\cos x = x\).}
\pair{numerical-calculus/findroot/2}{The same, to fifty digits---Newton's quadratic convergence makes
extra precision cheap.}
\pair{numerical-calculus/findroot/3}{A bracket \([0,1]\) selects Brent's method, which cannot leave the
bracket.}
\pair{numerical-calculus/findroot/4}{Two starting values select the derivative-free secant method.}
\pair{numerical-calculus/findroot/5}{A \(2\times2\) system solved by Newton with an exact Jacobian.}
\end{codepairs}

\usagebox{FindRoot}

\subsection{Optimization}
\label{ssec:optimization}

The last family finds minima and maxima. The distinction that organizes it is \emph{local} versus
\emph{global}. \B{FindMinimum}\index{FindMinimum@\mcode{FindMinimum}!methods} and \B{FindMaximum} start
from a point and descend to the nearest optimum; they are fast and precise but see only their own valley.
\B{NMinimize}\index{NMinimize@\mcode{NMinimize}!methods} and \B{NMaximize} search the whole region for the
\emph{global} optimum; they are slower but not fooled by local traps.

Start local. \B{FindMinimum}'s default is a quasi-Newton (BFGS) method in more than one variable, Brent's
method in one; it handles box and general constraints, and honors \texttt{WorkingPrecision}:

\begin{codepairs}
\pair{numerical-calculus/findminimum/1}{The Rosenbrock ``banana'' function, a standard hard case: the
minimum sits at \((1,1)\) in a curved valley, and BFGS threads it.}
\pair{numerical-calculus/findminimum/2}{A constrained problem: minimize \(x\cos x\) on \([1,15]\).}
\pair{numerical-calculus/findminimum/3}{Fifty digits of the minimizer of \((x-\pi)^2\)---it is \(\pi\).}
\pair{numerical-calculus/findminimum/4}{\B{FindMaximum} is the mirror image: the first peak of
\(\sin(x)\,e^{-x/10}\).}
\end{codepairs}

\usagebox{FindMinimum}

Beyond the default, \B{FindMinimum} offers a large menu of local methods---\texttt{"Newton"},
\texttt{"ConjugateGradient"}, \texttt{"LBFGSB"}, \texttt{"NelderMead"}, and the constrained solvers
\texttt{"COBYLA"}, \texttt{"COBYQA"}, and \texttt{"SLSQP"}, among others. A few names it recognizes
(\texttt{"LevenbergMarquardt"}, \texttt{"InteriorPoint"}) are not yet implemented and say so.

Now global. A function like Rastrigin's, whose graph is an egg-carton of local minima over a single global
one, would trap any purely local method; \B{NMinimize}'s default differential-evolution search is built to
escape:

\begin{codepairs}
\pair{numerical-calculus/nminimize/1}{The Rastrigin function in two variables has its global minimum
\(0\) at the origin, buried among hundreds of local minima; differential evolution finds it.}
\pair{numerical-calculus/nminimize/2}{A smooth constrained problem: minimize \(x+y\) on the unit disk,
answer \(-\sqrt2\) at \((-\tfrac1{\sqrt2},-\tfrac1{\sqrt2})\).}
\pair{numerical-calculus/nminimize/3}{Integer programming: the same objective restricted to the integer
lattice by \texttt{Element[\{x,y\}, Integers]}.}
\pair{numerical-calculus/nminimize/4}{The deterministic \texttt{"DIRECT"} method solves Rastrigin on a
box exactly, with no randomness at all.}
\end{codepairs}

\usagebox{NMinimize}

\B{NMinimize}'s global methods---\texttt{"DifferentialEvolution"}, \texttt{"SimulatedAnnealing"},
\texttt{"DualAnnealing"}, \texttt{"BasinHopping"}, \texttt{"SHGO"}, \texttt{"DIRECT"}, and others---each
take a bag of sub-options through \texttt{Method -> \{"Name", "SearchPoints" -> \dots\}}; the usage card
above and the reference page enumerate them.

\calloutnote{Theory}{Differential evolution\index{differential evolution} keeps a population of candidate
points and breeds new ones by adding the scaled difference of two random members to a third---a
gradient-free move that adapts its own step size to the population's spread~\cite{stornprice}. Simulated
annealing~\cite{kirkpatrick} instead walks one point randomly, accepting uphill moves with a probability
that cools over time, so it can climb out of a local basin early and settle late. And DIRECT~\cite{direct}
is deterministic: it repeatedly bisects the most promising rectangles of the domain, balancing global reach
against local refinement---which is why pair~4 lands exactly on the origin every time.}

\calloutnote{Performance}{Global optimizers evaluate the objective thousands of times, so Mathilda compiles
it to bytecode (Chapter~\ref{ch:compilation}) before the search begins. Supplying an \texttt{EvaluationMonitor}
forces every evaluation back through the interpreter and disables that fast path---useful for debugging a
search, costly in a production run.}

\subsection*{Where this connects}

Numerical calculus is the pragmatic mirror of the exact mathematics in the rest of this chapter. Each
function here has an exact sibling in \S\ref{sec:calculus}---\B{ND} to \B{D}, \B{NIntegrate} to
\B{Integrate}, \B{NLimit} to \B{Limit}, \B{NResidue} to \B{Residue}---and reaches for a number precisely
where the exact route runs out. Underneath, the numerical solvers lean on the machinery of the neighboring
sections: \B{NSolve} on the Gröbner bases of \S\ref{sec:algebra}, and \B{NRoots} and \B{NSolve} on the
eigenvalue kernels of \S\ref{sec:linalg}. For the exhaustive options and corner cases of any function named
here, its reference page is one \texttt{?Name} away at the prompt.
